% Generated by scripts/build_resume.py from _data/resume_en.yml; edit that file instead.
% Layout: Jake's Resume (MIT license). Compile with pdfLaTeX.
\documentclass[a4paper,11pt]{article}
\usepackage{latexsym}
\usepackage[empty]{fullpage}
\usepackage{titlesec}
\usepackage{marvosym}
\usepackage[usenames,dvipsnames]{color}
\usepackage{verbatim}
\usepackage{enumitem}
\usepackage[hidelinks]{hyperref}
\usepackage{fancyhdr}
\usepackage[english]{babel}
\usepackage{tabularx}
\input{glyphtounicode}

\pagestyle{fancy}
\fancyhf{}
\fancyfoot{}
\renewcommand{\headrulewidth}{0pt}
\renewcommand{\footrulewidth}{0pt}

\addtolength{\oddsidemargin}{-0.5in}
\addtolength{\evensidemargin}{-0.5in}
\addtolength{\textwidth}{1in}
\addtolength{\topmargin}{-.5in}
\addtolength{\textheight}{1.0in}

\urlstyle{same}
\raggedbottom
\raggedright
\setlength{\tabcolsep}{0in}

\titleformat{\section}{\vspace{-4pt}\scshape\raggedright\large}{}{0em}{}[\color{black}\titlerule \vspace{-5pt}]
\pdfgentounicode=1

\newcommand{\resumeItem}[1]{\item\small{#1 \vspace{-2pt}}}
\newcommand{\resumeSubheading}[4]{
  \vspace{-2pt}\item
    \begin{tabular*}{0.97\textwidth}[t]{l@{\extracolsep{\fill}}r}
      \textbf{#1} & #2 \\
      \textit{\small#3} & \textit{\small #4} \\
    \end{tabular*}\vspace{-7pt}
}
\newcommand{\resumeProjectHeading}[2]{
    \item
    \begin{tabular*}{0.97\textwidth}{l@{\extracolsep{\fill}}r}
      \small#1 & #2 \\
    \end{tabular*}\vspace{-7pt}
}
\newcommand{\resumeSubHeadingListStart}{\begin{itemize}[leftmargin=0.15in, label={}]}
\newcommand{\resumeSubHeadingListEnd}{\end{itemize}}
\newcommand{\resumeItemListStart}{\begin{itemize}}
\newcommand{\resumeItemListEnd}{\end{itemize}\vspace{-5pt}}

\begin{document}

\begin{center}
    \textbf{\Huge \scshape Bolin Li} \\ \vspace{1pt}
    \small bolin09511 at gmail dot com
\end{center}

\section{Education}
  \resumeSubHeadingListStart
    \resumeSubheading
      {University of Chinese Academy of Sciences}{Sep 2022 -- Jun 2025}
      {M.S., Computer Vision \& Computer Graphics}{}
    \resumeSubheading
      {Ocean University of China}{Sep 2018 -- Jun 2022}
      {B.S., Electronic Information}{}
  \resumeSubHeadingListEnd

\section{Experience}
  \resumeSubHeadingListStart
    \resumeSubheading
      {Guangzhou Tencent Technology Co., Ltd.}{Jun 2025 -- Present}
      {Game Engine Developer}{}
      \resumeItemListStart
        \resumeItem{\textbf{Engine features}: Customized Unreal Engine rendering features, including debug views, light shaft occlusion, translucent shadows, and light functions; ported engine plugins and adapted visual quality across platforms.}
        \resumeItem{\textbf{Performance, stability and tooling}: Performance and resources - skeletal mesh culling, shader permutation cleanup, offline shadow compression, HLOD unloading, texture downgrading, and spreading asset streaming across frames; separated per-platform assets and materials, managed consistently in the editor and at runtime. Tools and stability - built debug and statistics tools for triangle counts, pixel cost, and scene resource usage; diagnosed and fixed crashes related to numerical precision, asset import, shader translation, and graphics APIs, using RenderDoc, Unreal Insights, and ProfileGPU to analyze rendering anomalies and performance bottlenecks.}
      \resumeItemListEnd
    \resumeSubheading
      {Guangzhou Tencent Technology Co., Ltd.}{Mar 2024 -- Sep 2024}
      {Game Engine Intern - ported the mobile Overdraw debug view, enabled VRS}{}
  \resumeSubHeadingListEnd

\section{Projects}
  \resumeSubHeadingListStart
    \resumeProjectHeading
      {\textbf{Neural Material (Neural BSDF) Integration and Optimization in UE}}{Jun 2026 -- Present}
      \resumeItemListStart
        \resumeItem{\textbf{Integration and optimization}: Integrated neural materials into Unreal Engine and optimized their performance.}
        \resumeItem{\textbf{Benchmarking and archiving}: Fixed the rendering setup and used batched timing; sampled GPU clocks with nvidia-smi to flag downclocked runs; archived 190 reports covering 33 experiments at 2048px and 4 measurement protocols.}
        \resumeItem{\textbf{Comparability}: Found and handled two kinds of results that cannot be compared directly - different measurement protocols (about 109 ms at 1024px with per-frame timing versus about 17 ms at 2048px with batched timing) and session drift (an unchanged control material varied by about 15\%, between 7.7 and 8.8 ms); regression checks now use the cost ratio to the control material.}
        \resumeItem{\textbf{Regression checks}: Pinned a baseline with a 5\% tolerance check that returns a non-zero exit code on regression, so it can run in CI; an artificial 15\% slowdown was detected as +14.0\%.}
      \resumeItemListEnd
    \resumeProjectHeading
      {\textbf{AI-assisted UE Development and Troubleshooting System}}{Jul 2026 -- Present}
      \resumeItemListStart
        \resumeItem{\textbf{AI-assisted R\&D infrastructure}: Built AI-assisted development workflows for Unreal Engine (CodeGraph, UE Reviewer, UnrealMCP, TraceMCP, RenderTrail) covering code review, performance analysis, and issue diagnosis; explored closed-loop automated debugging and its adoption in the engine development process.}
        \resumeItem{\textbf{Closed loop across logic and visuals}: Runtime intermediate state cannot be fully recorded and log-based approaches require restarting the process. Used Rewind to record the final state of the faulty track, Checkpoints to save key states for iterative replay, and WinDbg to record the faulty flow, giving a semi-automated troubleshooting loop without stopping the game.}
        \resumeItem{\textbf{Structured rendering data}: Rendering bugs involve the RDG and GPU layers; combined RDG and frame-capture data, starting from the anomalous region to locate resource or shader problems, forming a troubleshooting loop on the rendering side.}
      \resumeItemListEnd
  \resumeSubHeadingListEnd

\end{document}
